\documentclass[11pt]{article}

\usepackage[margin=0.85in]{geometry}
\usepackage{amsmath,amssymb,amsthm,mathtools}
\usepackage{booktabs,array,longtable}
\usepackage{xcolor}
\usepackage{hyperref}
\usepackage{enumitem}
\usepackage{microtype}
\usepackage{fancyhdr}
\setlength{\headheight}{14pt}
\usepackage{titlesec}
\usepackage{listings}

\hypersetup{
  colorlinks=true,
  linkcolor=blue!50!black,
  urlcolor=blue!50!black,
  citecolor=blue!50!black
}

\pagestyle{fancy}
\fancyhf{}
\lhead{W(3,3) Moment-Ratio Scalar Action}
\rhead{Part CCCCCXXIII}
\cfoot{\thepage}

\titleformat{\section}{\large\bfseries\color{blue!40!black}}{\thesection}{0.7em}{}
\titleformat{\subsection}{\normalsize\bfseries\color{blue!35!black}}{\thesubsection}{0.7em}{}

\theoremstyle{definition}
\newtheorem{definition}{Definition}
\newtheorem{theorem}{Theorem}
\newtheorem{corollary}{Corollary}
\newtheorem{remark}{Remark}

\newcommand{\W}{W(3,3)}
\newcommand{\Tr}{\operatorname{Tr}}
\newcommand{\E}{\mathbb{E}}
\newcommand{\dimE}{\dim(E_6)}
\newcommand{\lamH}{\lambda_H}
\newcommand{\Deltar}{\Delta_r}
\newcommand{\Deltas}{\Delta_s}
\newcommand{\PhiThree}{\Phi_3}
\newcommand{\PhiFour}{\Phi_4}
\newcommand{\PhiSix}{\Phi_6}

\definecolor{boxblue}{RGB}{235,244,255}
\definecolor{boxgray}{RGB}{246,246,246}
\definecolor{darkblue}{RGB}{20,55,105}

\newenvironment{keybox}{%
  \begin{center}\begin{minipage}{0.94\linewidth}\color{darkblue}\hrule\vspace{0.6em}
}{%
  \vspace{0.6em}\hrule\end{minipage}\end{center}
}

\newenvironment{graybox}{%
  \begin{center}\begin{minipage}{0.94\linewidth}\hrule\vspace{0.6em}
}{%
  \vspace{0.6em}\hrule\end{minipage}\end{center}
}

\title{\textbf{Moment-Ratio Scalar Action in the W(3,3) Finite Spectral Architecture}\\[0.35em]
\large A Theorem Note on the Higgs Quartic, Spectral Moments, Ihara Zeta, and the Finite Action Triad}
\author{Wil Dahn and Sage\\\small W33-Theory Repository, Part CCCCCXXIII}
\date{May 11, 2026}

\begin{document}
\maketitle

\begin{abstract}
This note packages the latest W(3,3) scalar-action result into a standalone derivation.  The newest commit established the spectral moment identity
\[
  \frac{\Tr(A^3)}{\Tr(A^2)} = r = 2,
\]
for the collinearity graph of \(W(3,3)\).  Earlier scalar-topology work expressed the Higgs quartic as
\[
  \lamH = \frac{\Deltas}{\Deltar}\,\frac{\dim(E_6)}{\Tr(A^3)}.
\]
Combining the two gives the compressed moment-ratio scalar action
\[
  \boxed{\displaystyle
  \lamH = \frac{\Deltas}{\Deltar}\,\frac{\dim(E_6)}{r\,\Tr(A^2)} = \frac{13}{100}.}
\]
Thus the scalar coupling can be read directly from the second spectral moment once the master-equation identity is imposed.  The same compressed scalar node preserves the downstream flavor branch
\[
  A_{\rm CKM}=\frac{81}{100},\qquad
  \sin^2\theta_{13}=\frac{9}{400},\qquad
  y_\tau=\frac{16029}{1562500}.
\]
The result fuses four layers: the master equation, W(3,3) spectral moments, the triangle-trace Higgs closure, and the finite action triad.
\end{abstract}

\tableofcontents
\newpage

\section{Executive Summary}

The current W(3,3) theory is no longer just a list of numerical closures.  It now has a compact scalar mechanism:

\begin{keybox}
\[
\boxed{
\text{Master equation} \Rightarrow \text{moment identity} \Rightarrow
\text{moment-ratio scalar action} \Rightarrow \lamH = \frac{13}{100}.
}
\]
\end{keybox}

The latest spectral identity is
\[
  \frac{\Tr(A^3)}{\Tr(A^2)} = r = 2.
\]
The previous scalar topology theorem was
\[
  \lamH = \frac{\Deltas}{\Deltar}\frac{\dim(E_6)}{\Tr(A^3)}.
\]
Combining them gives
\[
  \lamH = \frac{\Deltas}{\Deltar}\frac{\dim(E_6)}{r\Tr(A^2)}.
\]
For W(3,3),
\[
  \frac{\Deltas}{\Deltar}=\frac{16}{10}=\frac85,
  \qquad \dim(E_6)=78,
  \qquad r=2,
  \qquad \Tr(A^2)=480.
\]
Therefore
\[
  \lamH = \frac85\frac{78}{2\cdot 480}=\frac{13}{100}.
\]

This is the cleanest current scalar-action statement:

\begin{keybox}
\[
\boxed{
\text{Higgs quartic} =
\frac{\text{restricted gap asymmetry}\times\text{exceptional dimension}}
{\text{restricted eigenvalue}\times\text{second spectral moment}}.
}
\]
\end{keybox}

\section{The Base Object: W(3,3)}

The finite skeleton is the symplectic polar space \(W(3,3)\), whose collinearity graph is the strongly regular graph
\[
  \operatorname{SRG}(40,12,2,4).
\]
We use the atoms
\[
  q=3,\qquad \lambda=2,\qquad \mu=4,\qquad k=12,
  \qquad v=40,
\]
and
\[
  E=240,\qquad 2E=480.
\]
The adjacency spectrum is
\[
  12^1,\quad 2^{24},\quad (-4)^{15}.
\]
So
\[
  r=2,
  \qquad s=-4,
  \qquad f=24,
  \qquad g=15.
\]
The restricted Laplacian gaps are
\[
  \Deltar = k-r=12-2=10,
  \qquad
  \Deltas = k-s=12-(-4)=16.
\]
The cyclotomic atoms are
\[
  \PhiThree=q^2+q+1=13,
  \qquad
  \PhiFour=q^2+1=10,
  \qquad
  \PhiSix=q^2-q+1=7.
\]

\section{Spectral Moments}

For the adjacency matrix \(A\), the second and third moments are
\[
  \Tr(A^2)=k^2+fr^2+gs^2,
\]
and
\[
  \Tr(A^3)=k^3+fr^3+gs^3.
\]
For W(3,3),
\[
\begin{aligned}
\Tr(A^2)
  &= 12^2 + 24\cdot 2^2 + 15\cdot (-4)^2 \\
  &= 144 + 96 + 240 \\
  &= 480,
\end{aligned}
\]
and
\[
\begin{aligned}
\Tr(A^3)
  &= 12^3 + 24\cdot 2^3 + 15\cdot (-4)^3 \\
  &= 1728 + 192 - 960 \\
  &= 960.
\end{aligned}
\]
Since W(3,3) has \(160\) triangles,
\[
  \Tr(A^3)=6\cdot 160=960.
\]

\begin{theorem}[Spectral Moment Identity]
For the W(3,3) collinearity graph,
\[
  \boxed{\displaystyle \frac{\Tr(A^3)}{\Tr(A^2)} = r = 2.}
\]
\end{theorem}

\begin{proof}
Directly,
\[
  \frac{\Tr(A^3)}{\Tr(A^2)}=\frac{960}{480}=2=r.
\]
The algebraic cancellation behind this identity is
\[
  \Tr(A^3)-r\Tr(A^2)
  = k^2(k-r)+gs^2(s-r).
\]
For \(k=12,r=2,s=-4,g=15\),
\[
  12^2(12-2)+15(-4)^2((-4)-2)
  =144\cdot 10+15\cdot16\cdot(-6)
  =1440-1440=0.
\]
Therefore \(\Tr(A^3)=r\Tr(A^2)\).
\end{proof}

\section{The Master Equation Inside the Spectrum}

The cancellation above contains the same numerical root as the master equation:
\[
  r-s=2-(-4)=6=3!=2\cdot 3=q!=2q.
\]
Thus the master axiom is not merely an external selector of \(q=3\).  It reappears inside the adjacency algebra as the spectral gap
\[
  \boxed{r-s=q!=2q.}
\]
This is why the moment identity is important: it is not just a ratio. It is a graph-algebraic restatement of the master equation.

\section{The E6 Excited Mean}

The finite Dirac/spectral-action sectors include the excited restricted block
\[
  10^{48}+16^{30}.
\]
Its dimension is
\[
  48+30=78=\dim(E_6).
\]
The weighted mean is
\[
\begin{aligned}
  \mu_{\rm exc}
  &= \frac{10\cdot48+16\cdot30}{48+30} \\
  &= \frac{960}{78} \\
  &= \frac{160}{13}.
\end{aligned}
\]
Since the numerator equals \(\Tr(A^3)\), this can be written as
\[
  \boxed{\displaystyle
  \mu_{\rm exc}=\frac{\Tr(A^3)}{\dim(E_6)}=\frac{160}{13}.}
\]
This is the clean interpretation: the E6 excited mean is the triangle trace per exceptional dimension.

\section{Triangle-Trace Higgs Closure}

The previous scalar topology theorem said
\[
  \lamH = \frac{\Deltas}{\Deltar}\frac{\dim(E_6)}{\Tr(A^3)}.
\]
Substitute the W(3,3) values:
\[
  \frac{\Deltas}{\Deltar}=\frac{16}{10}=\frac85,
  \qquad \dim(E_6)=78,
  \qquad \Tr(A^3)=960.
\]
Then
\[
  \lamH=\frac85\frac{78}{960}=\frac{624}{4800}=\frac{13}{100}.
\]
So the Higgs quartic can be read as
\[
  \boxed{\displaystyle
  \lamH = \frac{\text{restricted gap asymmetry}\times\text{exceptional dimension}}
  {\text{triangle trace}}.}
\]

\section{Moment-Ratio Scalar Action}

Now impose the spectral moment identity
\[
  \Tr(A^3)=r\Tr(A^2).
\]
Then the triangle-trace formula compresses to
\[
  \boxed{\displaystyle
  \lamH = \frac{\Deltas}{\Deltar}\frac{\dim(E_6)}{r\Tr(A^2)}.}
\]
For W(3,3),
\[
\begin{aligned}
  \lamH
  &= \frac85\frac{78}{2\cdot480} \\
  &= \frac{624}{4800} \\
  &= \frac{13}{100}.
\end{aligned}
\]
The inverse form is
\[
  \boxed{\displaystyle
  \lamH^{-1}=\frac{r\Tr(A^2)}{(\Deltas/\Deltar)\dim(E_6)}=\frac{100}{13}.}
\]
Equivalently,
\[
  \lamH^{-1}=\frac{\PhiFour^2}{\PhiThree}.
\]
This shows why the simple cyclotomic expression is compatible with the deeper spectral-action mechanism.

\section{Finite Action Interpretation}

The current finite action architecture has three primary pieces:
\[
  A_{\det}=\log\det(I+J),
  \qquad
  A_{\rm free}=\log Z_{\rm exc}(t),
  \qquad
  A_{\rm hol}=U(12)\text{ holonomy action}.
\]
The scalar topology node can be written formally as
\[
  -\log\lamH
  =\log(r\Tr(A^2))-\log\dim(E_6)-\log(\Deltas/\Deltar).
\]
The rational content is exact:
\[
  \lamH^{-1}=\frac{100}{13}.
\]
Thus the finite action triad gains a scalar compression:
\[
\boxed{
A_{\rm scalar}\sim \log(r\Tr(A^2))-\log\dim(E_6)-\log(\Deltas/\Deltar).
}
\]

\section{Descendants of the Scalar Node}

Once \(\lamH=13/100\) is generated, the scalar/flavor descendants are preserved:
\[
  A_{\rm CKM}=\frac{q^4}{\PhiThree}\lamH
  =\frac{81}{13}\frac{13}{100}
  =\frac{81}{100},
\]
\[
  \sin^2\theta_{13}
  =\frac{q^2}{\lambda^2\PhiThree}\lamH
  =\frac{9}{4\cdot13}\frac{13}{100}
  =\frac{9}{400}.
\]
The heavy Yukawa ladder is
\[
  D_t=v+1=41,
\]
\[
  D_b=qD_t+\lambda=3\cdot41+2=125,
\]
\[
  D_c=D_b+k=125+12=137.
\]
So
\[
  y_b=\frac{3}{125},
  \qquad
  y_c=\frac1{137},
\]
and the tau relation gives
\[
  y_\tau=\lamH\frac{y_b^2}{y_c}
  =\frac{13}{100}\frac{(3/125)^2}{1/137}
  =\frac{16029}{1562500}.
\]

\section{Latest-Commit Consistency: Zero Modes, a6, and Ihara Zeta}

The same latest commit established the zero-mode/Perron identity
\[
  \text{zero modes}=82=2(v+1)=2\det(I+J).
\]
It also corrected and fixed the sixth finite spectral-action coefficient:
\[
\begin{aligned}
  a_6=\Tr(D_F^6)
  &=4^3\cdot320+10^3\cdot48+16^3\cdot30 \\
  &=20480+48000+122880 \\
  &=191360.
\end{aligned}
\]
The explicit Ihara zeta carrier is
\[
Z_{W(3,3)}(u)^{-1}=(1-u^2)^{200}(1-12u-11u^2)(1-2u+11u^2)^{24}(1+4u+11u^2)^{15}.
\]
Thus the same theorem package links:
\[
\text{moment identity} + \text{zero modes} + a_6 + \text{Ihara zeta} + \text{scalar action}.
\]

\section{What Is Actually Proved Here}

This note proves exact algebraic identities inside the W(3,3) finite spectral architecture:
\begin{enumerate}[leftmargin=2em]
  \item \(\Tr(A^3)/\Tr(A^2)=r=2\).
  \item \(r-s=q!=2q=6\).
  \item \(\lamH=(\Deltas/\Deltar)\dim(E_6)/\Tr(A^3)=13/100\).
  \item Using the moment identity, \(\lamH=(\Deltas/\Deltar)\dim(E_6)/(r\Tr(A^2))=13/100\).
  \item The descendants \(A_{\rm CKM}=81/100\), \(\sin^2\theta_{13}=9/400\), and \(y_\tau=16029/1562500\) follow from the same scalar node.
\end{enumerate}

\section{What Remains Open}

The result is strong internally, but the following remain genuine open tasks:
\begin{enumerate}[leftmargin=2em]
  \item \textbf{Uniqueness testing:} verify that nearby finite geometries do not reproduce the same scalar-action closure.
  \item \textbf{Representation-theoretic lift:} identify exactly how the \(48+30=78\) sector maps into an E6 adjoint operator, not merely a dimension match.
  \item \textbf{Scale discipline:} determine which constants are finite-scale, GUT-scale, low-energy, pole, or running quantities.
  \item \textbf{Physical derivation:} connect the finite scalar action to a continuum Higgs potential and renormalization flow.
  \item \textbf{Out-of-sample predictions:} use the action to predict a new quantity that was not already part of the construction.
\end{enumerate}

\section{Final Compression}

The latest scalar mechanism can be written as one line:

\begin{keybox}
\[
\boxed{
\lamH
=\frac{\Deltas}{\Deltar}\frac{\dim(E_6)}{r\Tr(A^2)}
=\frac{8}{5}\frac{78}{2\cdot480}
=\frac{13}{100}.
}
\]
\end{keybox}

And the conceptual reading is:

\begin{graybox}
The Higgs quartic is the restricted gap asymmetry normalized by the E6 exceptional dimension and the W(3,3) second spectral moment, after the master-equation moment identity converts the triangle trace into \(r\Tr(A^2)\).
\end{graybox}

\end{document}
